\documentclass[11pt]{article}
\usepackage[margin=1in]{geometry}
\usepackage[T1]{fontenc}
\usepackage{lmodern,microtype,amsmath,amssymb,amsthm,mathtools,booktabs,array}
\usepackage[hidelinks]{hyperref}
\hypersetup{pdftitle={RA-01: Tail-envelope guarantees and a clock obstruction},pdfsubject={Partial results; not a complete resolution of RA-01}}
\usepackage{enumitem}
\setlength{\parindent}{0pt}
\setlength{\parskip}{5pt}
\setlength{\emergencystretch}{2em}
\newtheorem{theorem}{Theorem}[section]
\newtheorem{lemma}[theorem]{Lemma}
\newtheorem{proposition}[theorem]{Proposition}
\newtheorem{corollary}[theorem]{Corollary}
\theoremstyle{remark}\newtheorem{remark}[theorem]{Remark}
\DeclareMathOperator{\tr}{tr}
\DeclareMathOperator{\rank}{rank}
\newcommand{\E}{\mathbb E}
\newcommand{\Prob}{\mathbb P}
\newcommand{\C}{\mathbb C}
\newcommand{\R}{\mathbb R}
\newcommand{\eps}{\varepsilon}
\newcommand{\Id}{I}
\title{RA-01: Tail-envelope guarantees and a clock obstruction}
\author{}\date{}
\begin{document}
\begin{center}
{\Large\bfseries RA-01: Tail-envelope guarantees and a clock obstruction}
\end{center}
\vspace{0.4em}
\begin{abstract}
The unrestricted RPCholesky optimal-pivot-count question is not resolved in this note. We establish the requested expected relative-error guarantee with the explicit constant $C=3$ when the eigenvalues after the target rank lie below a half-geometric envelope, and with $C=4$ under an inverse-square tail envelope. The leading eigenvalues may have arbitrarily large spread, and all eigenvectors are unrestricted. A further theorem gives a rank-independent constant for every fixed summable power envelope, with its dependence on the envelope exponent displayed explicitly. These results retain a head-product denominator in the determinant-path argument rather than imposing decay on the entire spectrum. We also construct finite rank-one-spike examples disproving every affine expected-count estimate with unit tail-rate coefficient, even when its additive rank coefficient is arbitrarily large. A saturation calculation explains why the determinant comparison used here does not settle the diffuse-tail case. Exact rational checks and real/complex floating-point diagnostics accompany the proofs. The general arguments have not been independently reviewed or proof-assistant verified.
\end{abstract}

\section{Target, scope, and relation to preceding work}
Let $A\in\C^{n\times n}$ be Hermitian positive semidefinite, with eigenvalues
$\lambda_1\ge\cdots\ge\lambda_n\ge0$. Define
\[
 \tau_r(A)=\sum_{j>r}\lambda_j(A),\qquad \tau_0(A)=\tr A.
\]
RPCholesky starts with $R_0=A$. At a nonzero residual $R$, set $T=\tr R$, choose a coordinate $j$ with probability $R_{jj}/T$, and update
\begin{equation}
 R^{(j)}=R-\frac{Re_je_j^*R}{R_{jj}}. \label{eq:update}
\end{equation}
A zero residual stays zero. Write $T_k=\tr R_k$ and $f_k(A)=\E T_k$.
The canonical RA-01 question asks for one absolute $C\ge1$ such that
\begin{equation}
 f_{\min\{n,\lceil Cr/\eps\rceil\}}(A)
 \le (1+\eps)\tau_r(A)
 \quad(1\le r\le n,\ 0<\eps<1) \label{eq:target}
\end{equation}
for every allowed matrix, with exactly the law in \eqref{eq:update}. The repository entry inspected on September 13, 2026 displays \emph{Open}.\footnote{Canonical statement and status: \cite{repo}.}

\begin{center}
\begin{tabular}{@{}p{0.51\linewidth}p{0.42\linewidth}@{}}
\toprule
Scope of the result & Conclusion established here\\
\midrule
$\lambda_{r+j}\le\lambda_{r+1}2^{-(j-1)}$ & \eqref{eq:target} with $C=3$.\\[3pt]
$\lambda_{r+j}\le\lambda_{r+1}j^{-2}$ & \eqref{eq:target} with $C=4$.\\[3pt]
$\lambda_{r+j}\le\lambda_{r+1}j^{-p}$, fixed $p>1$ & An explicit finite constant $C_p$, independent of rank, dimension, and head spread.\\[3pt]
All matrices; continuous-time pivot count & No finite $a$ gives $\E N(t)\le ar+t\tau_r(A)$ universally.\\[3pt]
All matrices; unrestricted RA-01 & No universal upper constant or counterexample to all constants is proved.\\
\bottomrule
\end{tabular}
\end{center}

The determinant-path framework is due to the cited RPCholesky analyses. Epperly's August 2026 result gives a sufficient count containing an additional $2r\sqrt{\log r}$ term.\footnote{See Corollary 1.3 and the conjecture after (1.5) in \cite{epperly}.} Divan and Gilles give elementary-symmetric-polynomial estimates and convergence rates under decay of the full spectrum.\footnote{See Theorem 3.1, Lemma 3.2, Corollary 3.3, and Appendix A in \cite{divan}.} Here decay is assumed only \emph{after} rank $r$, and the scale is the first tail eigenvalue. The resulting relative-error bounds do not depend on the first $r$ eigenvalues.

The preceding packages contain, among other results, a necessary lower bound $C\ge2$, a bounded-leading-spread upper theorem, and a coordinate-block clock argument. They remain unpublished, not independently reviewed arguments. They are retained for provenance; the present note does not independently re-audit their all-subsets lower bound. No literature-priority claim is made for the results or methods here.

\section{Residual identities and scalar completion}
\begin{lemma}[Order, rank, and spectral tail]
Along every realized path, $0\preceq R_{k+1}\preceq R_k\preceq A$. A positive pivot reduces the rank by exactly one. For $0\le k\le n$,
\begin{equation}
 T_k\ge\tau_k(A),\qquad \tau_r(R_k)\le\tau_r(A). \label{eq:tailorder}
\end{equation}
In particular, $R_k=0$ for $k\ge\rank A$.
\end{lemma}
\begin{proof}
For a positive pivot, put $u=R^{1/2}e_j/\sqrt{R_{jj}}$. Then $u$ is a unit vector in the range of $R$, and $R^{(j)}=R^{1/2}(I-uu^*)R^{1/2}$. This proves the order and rank assertions. The matrix $A-R_k$ is PSD with rank at most $k$ and is bounded above by $A$. Its trace is at most the sum of the first $k$ eigenvalues of $A$, giving the first inequality in \eqref{eq:tailorder}. Eigenvalue monotonicity under PSD order gives the second.
\end{proof}

\begin{lemma}[Scalar completion]\label{lem:scalar}
Fix $r\ge1$ and $\tau=\tau_r(A)>0$. If $f_{t_0}\le B\tau$ for $B>1$, then, for every integer $s\ge0$,
\begin{equation}
 f_{t_0+s}\le\tau\left(1+
 \frac{1}{(B-1)^{-1}+s/(rB)}\right). \label{eq:completion}
\end{equation}
It suffices to take $s\ge rB(\eps^{-1}-(B-1)^{-1})$ when this quantity is positive. The conservative condition $s\ge rB/\eps$ always suffices.
\end{lemma}
\begin{proof}
For a nonzero residual of trace $T$, averaging \eqref{eq:update} gives
\begin{equation}
 \E[T_{k+1}\mid R_k=R]=T-\frac{\tr R^2}{T}. \label{eq:onedrift}
\end{equation}
Zero diagonal entries of a PSD matrix have zero columns, so omitting them loses no term. From \eqref{eq:tailorder}, $\tr R^2\ge (T-\tau)_+^2/r$. The function
\[
 g(x)=\begin{cases}
 x,&0\le x\le\tau,\\
 x-(x-\tau)^2/(rx),&x>\tau
 \end{cases}
\]
is nondecreasing and concave. On the second branch it equals $(1-1/r)x+2\tau/r-\tau^2/(rx)$, and the first derivatives agree at $\tau$. Jensen's inequality yields $f_{k+1}\le g(f_k)$.

Set $u_0=B-1$ and $u_{s+1}=u_s-u_s^2/[r(1+u_s)]$. These numbers are positive and nonincreasing, and monotonicity of $g$ yields $f_{t_0+s}\le\tau(1+u_s)$. Direct calculation gives
\[
 \frac1{u_{s+1}}-\frac1{u_s}
 =\frac1{r+(r-1)u_s}\ge\frac1{rB}.
\]
Summation proves \eqref{eq:completion}.
\end{proof}

\begin{lemma}[Two small-rank starting bounds]\label{lem:small}
For every PSD matrix,
\begin{equation}
 f_k\le(1+1/k)\tau_1(A)\quad(k\ge1),
 \qquad f_7\le\frac{89}{56}\tau_2(A)<\frac85\tau_2(A). \label{eq:small}
\end{equation}
\end{lemma}
\begin{proof}
For rank one, let $\lambda=\lambda_1(A)$ and $\tau=\tau_1(A)>0$. Equation \eqref{eq:onedrift} gives $f_1\le\lambda+\tau-\lambda^2/(\lambda+\tau)\le2\tau$. For rank one the scalar map sends $\tau(1+1/k)$ to $\tau(1+1/(k+1))$. Induction proves the first assertion; zero tails are settled by termination.

For completeness, the error-doubling argument needed for rank two can be checked directly.\footnote{This is the mechanism of the established error-doubling lemma, \cite[Lemma 5.5]{chen}.} Let $P$ project onto the first $q$ eigenvectors of a PSD matrix $R$, write $R=H+E$ for its spectral head and tail, and put $\tau=\tr E$, $T=\tr R$. For a pivot $j$ with $c_j=P_{jj}>0$, use $Pe_j/\sqrt{c_j}$ together with an orthonormal basis of $\operatorname{range}(I-P)$ as trial vectors in the variational formula for $\tau_{q-1}(R^{(j)})$. If $h_j=H_{jj}$ and $b_j=E_{jj}$, this gives
\[
 \tau_{q-1}(R^{(j)})\le\tau+\frac{h_jb_j}{c_j R_{jj}}.
\]
Since $h_j/c_j\le\lambda_1(R)$, the expected second term is at most $\lambda_1(R)\tau/T$. When $c_j=0$, use a leading eigenvector instead; its contribution has the same bound. Thus $\E\tau_{q-1}(R^{(j)})\le2\tau_q(R)$. Applying this with $q=2$ and then the rank-one one-pivot bound gives $f_2\le4\tau_2(A)$.

Normalize $\tau_2(A)=1$. The rank-two scalar map gives the following valid upper comparisons, with deliberate round-ups at steps four and five:
\[
\begin{array}{c|rrrrrr}
 k&2&3&4&5&6&7\\\hline
 f_k\text{ upper bound}&4&23/8&16/7&2&7/4&89/56.
\end{array}
\]
Indeed $g(23/8)=833/368<16/7$ and $g(16/7)=431/224<2$. This proves the second assertion.
\end{proof}

\section{A determinant-path bound retaining the head denominator}
Write $e_j(A)$ for the $j$th elementary symmetric polynomial of the eigenvalues, with $e_0=1$. It is also the sum of all principal minors of order $j$. This follows by expanding $\det(I+zA)$ in two ways.

\begin{proposition}[Path probability and integrated bound]\label{prop:path}
Fix $\tau=\tau_r(A)>0$, $k=r+m\le n$, and $m\ge1$. Put
\begin{equation}
 D_r(A)=\prod_{i=0}^{r-1}\tau_i(A),\qquad
 M_{r,m}(A)=\frac{k!e_k(A)}{D_r(A)\tau^m}. \label{eq:M}
\end{equation}
Then, for $v>0$,
\begin{equation}
 \Prob\{T_k>v\tau\}\le M_{r,m}(A)v^{-m}. \label{eq:tailprob}
\end{equation}
For $m>1$ this implies
\begin{equation}
 \frac{f_k}{\tau}\le\frac{m}{m-1}M_{r,m}(A)^{1/m}. \label{eq:integrated}
\end{equation}
\end{proposition}
\begin{proof}
A valid ordered path $j_1,\ldots,j_k$, with unordered set $S$, has probability
\[
 \prod_{i=0}^{k-1}\frac{(R_i)_{j_{i+1}j_{i+1}}}{T_i}
 =\frac{\det A_{SS}}{\prod_{i=0}^{k-1}T_i}.
\]
The numerator is the product of the Cholesky pivots. On the event $T_k>v\tau$, the first $r$ denominators obey $T_i\ge\tau_i(A)$ and the remaining $m$ obey $T_i\ge T_k>v\tau$. Sum the resulting upper bound over ordered paths. Each $k$-set has at most $k!$ orders, and the sum of its minors is $e_k(A)$, proving \eqref{eq:tailprob}. Terminated paths contribute no bad event.

For $m>1$, integrate $\min\{1,Mv^{-m}\}$ over $v>0$ and split the integral at $M^{1/m}$. The result is $\frac{m}{m-1}M^{1/m}$. If $M=0$, the error is zero.
\end{proof}

Let $a=\lambda_{r+1}(A)>0$, $P_r=\prod_{i=1}^r\lambda_i(A)$, and let $e_j^{\rm tail}$ denote the symmetric polynomial in the tail eigenvalues. Since $\lambda_i\ge a$ for $i\le r$,
\begin{equation}
 e_{r-s}(\lambda_1,\ldots,\lambda_r)
 \le {r\choose s}P_r a^{-s},\qquad D_r(A)\ge P_r. \label{eq:head}
\end{equation}
The expansion
\begin{equation}
 e_{r+m}(A)=\sum_{s=0}^r
 e_{r-s}(\lambda_1,\ldots,\lambda_r)e_{m+s}^{\rm tail} \label{eq:split}
\end{equation}
therefore eliminates all head magnitudes from \eqref{eq:M} once the tail coefficients are controlled. This is a spectral estimate on path probabilities in the original coordinates, not a claim that the pivot law is unitarily invariant.

\section{A half-geometric tail gives the constant three}
\begin{theorem}[Geometric-tail guarantee]\label{thm:geo}
Suppose $\tau_r(A)>0$, put $a=\lambda_{r+1}(A)$, and assume
\begin{equation}
 \lambda_{r+j}(A)\le a\,2^{-(j-1)}
 \quad(1\le j\le n-r). \label{eq:geoassumption}
\end{equation}
Then the original RPCholesky algorithm satisfies \eqref{eq:target} with $C=3$. The first $r$ eigenvalues and all eigenvectors are unrestricted. Zero-tail inputs satisfy the same conclusion by termination.
\end{theorem}

\begin{lemma}[Geometric coefficients]\label{lem:geocoeff}
Under \eqref{eq:geoassumption}, for every integer $\ell\ge1$,
\begin{equation}
 e_\ell^{\rm tail}\le
 \frac{a^\ell 2^{-\ell(\ell-1)/2}}{\prod_{j=1}^\ell(1-2^{-j})}
 <4a^\ell2^{-\ell(\ell-1)/2}. \label{eq:geocoeff}
\end{equation}
\end{lemma}
\begin{proof}
By coefficientwise monotonicity it suffices to use the infinite sequence $a, a/2,a/4,\ldots$. In the defining sum, write increasing indices as $j_i=i+b_i$, where $0\le b_1\le\cdots\le b_\ell$. Summing over the successive nonnegative gaps between the $b_i$ gives the product of geometric series $\prod_{j=1}^\ell(1-2^{-j})^{-1}$, with the fixed factor $2^{-\ell(\ell-1)/2}$.

The full infinite product has the lower bound
\[
 \prod_{j=1}^\infty(1-2^{-j})
 \ge\frac{21}{64}\left(1-\sum_{j=4}^\infty2^{-j}\right)
 =\frac{147}{512}>\frac14.
\]
Here the inequality $\prod(1-x_j)\ge1-\sum x_j$ follows first for finite products by induction and then by a limit. Every finite product is at least the infinite product.
\end{proof}

Combining \eqref{eq:head}--\eqref{eq:geocoeff} gives
\begin{align}
 e_{r+m}(A)
 &\le4P_r a^m2^{-m(m-1)/2}
 \sum_{s=0}^r{r\choose s}2^{-ms-s(s-1)/2}\notag\\
 &\le4P_r a^m2^{-m(m-1)/2}(1+2^{-m})^r. \label{eq:geoek}
\end{align}
Since $a\le\tau_r(A)$, Proposition~\ref{prop:path} now yields
\begin{equation}
 \frac{f_{r+m}}{\tau_r(A)}
 \le\frac{m}{m-1}
 \bigl[4(r+m)!(1+2^{-m})^r\bigr]^{1/m}
 2^{-(m-1)/2}\quad(m>1). \label{eq:geomean}
\end{equation}

\begin{lemma}[Uniform geometric-tail starting point]\label{lem:geowarm}
Under \eqref{eq:geoassumption},
\begin{equation}
 f_{3r+1}\le\frac85\tau_r(A). \label{eq:geowarm}
\end{equation}
\end{lemma}
\begin{proof}
Termination handles $n\le3r+1$. Lemma~\ref{lem:small} handles $r=1,2$. For $r\ge3$, take $m=2r+1$ in \eqref{eq:geomean}. Since $r2^{-(2r+1)}\le1/8$,
\[
 (1+2^{-(2r+1)})^r\le\frac1{1-r2^{-(2r+1)}}\le\frac87,
\]
using ${r\choose j}\le r^j$ and the infinite geometric series. Hence
\[
 \frac{f_{3r+1}}{\tau_r(A)}
 \le\frac{2r+1}{2r}[5(3r+1)!]^{1/(2r+1)}2^{-r}.
\]
Put $F_r=(8/5)(2r/(2r+1))2^r$. It is enough to prove
$5(3r+1)!\le F_r^{2r+1}$ for $r\ge3$. The exact base ratio is
\[
 \frac{5\,10!}{(384/35)^7}
 =\frac{56\,296\,884\,765\,625}{59\,373\,627\,899\,904}<1.
\]
Also $F_r\ge2^r$ and $F_{r+1}/F_r>2$. Thus the ratio between successive right sides is greater than $2^{4r+3}$. The factorial ratio on the left is at most $(3r+4)^3\le2^{3r+3}$ for $r\ge3$: the inequality $3r+4\le2^{r+1}$ holds at $r=3$ and is preserved by induction. This proves the required factorial inequality.
\end{proof}

\begin{proof}[Proof of Theorem~\ref{thm:geo}]
Let $k=\lceil3r/\eps\rceil$. The dimension cap and zero tail are immediate. Otherwise $k<n$, so $n>3r+1$, and $k\ge3r+1$ because $\eps<1$. Use Lemma~\ref{lem:scalar} with $t_0=3r+1$, $B=8/5$. The required extra count, when positive, is
\[
 \frac{8r}{5\eps}-\frac{8r}{3}.
\]
The available count exceeds it, since
\[
 k-(3r+1)-\left(\frac{8r}{5\eps}-\frac{8r}{3}\right)
 \ge\frac{7r}{5\eps}-\frac r3-1
 >\frac{16r}{15}-1>0.
\]
This proves the stated relative error with the exact prescribed count.
\end{proof}

\section{An inverse-square tail gives the constant four}
\begin{theorem}[Inverse-square-tail guarantee]\label{thm:square}
Suppose $\tau_r(A)>0$ and
\begin{equation}
 \lambda_{r+j}(A)\le a j^{-2},\qquad
 a=\lambda_{r+1}(A),\quad1\le j\le n-r. \label{eq:squareassumption}
\end{equation}
Then \eqref{eq:target} holds with $C=4$, with unrestricted leading eigenvalues and eigenvectors.
\end{theorem}
\begin{proof}
The classical infinite product for $\sinh$ gives\footnote{NIST DLMF, equation 4.36.1, \cite{dlmf}.}
\[
 \prod_{j=1}^\infty(1+z/j^2)
 =\frac{\sinh(\pi\sqrt z)}{\pi\sqrt z}
 =\sum_{\ell=0}^\infty\frac{\pi^{2\ell}z^\ell}{(2\ell+1)!}.
\]
The product converges uniformly on compact sets; its coefficients are nonnegative. Thus, under \eqref{eq:squareassumption},
\[
 e_\ell^{\rm tail}\le\frac{a^\ell\pi^{2\ell}}{(2\ell+1)!}
 \le\frac{a^\ell10^\ell}{(2\ell+1)!}.
\]
The elementary inequality $\pi<22/7$ suffices for $\pi^2<10$; for example it follows from the positive integral
$\int_0^1 x^4(1-x)^4/(1+x^2)\,dx=22/7-\pi$.

Using \eqref{eq:split} and the fact that each of the $2s$ additional factorial factors is at least $2m+2$ gives
\[
 e_{r+m}(A)\le P_r a^m\frac{10^m}{(2m+1)!}
 \left(1+\frac{10}{(2m+2)^2}\right)^r.
\]
Consequently, for $m>1$,
\begin{equation}
 \frac{f_{r+m}}{\tau_r(A)}\le
 \frac{10m}{m-1}
 \left[\frac{(r+m)!}{(2m+1)!}
 \left(1+\frac{10}{(2m+2)^2}\right)^r\right]^{1/m}. \label{eq:squaremean}
\end{equation}

We claim $f_{3r+1}\le2\tau_r(A)$. As before, termination and Lemma~\ref{lem:small} handle the capped case and ranks one and two. For $r\ge3$, take $m=2r+1$ in \eqref{eq:squaremean}. Since $10r/(4r+4)^2\le5/32$, the same geometric-series estimate gives
\[
 \left(1+\frac{10}{(4r+4)^2}\right)^r\le\frac{32}{27}<\frac65.
\]
Set $b_r=2r/[5(2r+1)]$ and $L_r=6(3r+1)!/[5(4r+3)!]$. The desired bound is equivalent to $L_r\le b_r^{2r+1}$. At $r=3$ the exact ratio is
\[
 \frac{L_3}{b_3^7}=\frac{367\,653\,125}{480\,370\,176}<1.
\]
The sequence $b_r$ increases, and
\[
 \frac{L_{r+1}}{L_r}
 =\frac{(3r+2)(3r+3)(3r+4)}{(4r+4)(4r+5)(4r+6)(4r+7)}.
\]
At $r=3$ this is $143/7752<(6/35)^2$. For $r\ge4$, use $3r+4\le13(r+1)/4$ to bound the ratio by
\[
 \frac{2197}{16384(r+1)}\le\frac{2197}{81920}
 <\frac{36}{1225}=b_3^2\le b_r^2.
\]
The ratio between successive $b_r^{2r+1}$ is greater than $b_r^2$. Induction proves the claim.

Finally set $k=\lceil4r/\eps\rceil$, with an inactive dimension cap. The warm start is at $t_0=3r+1$ with $B=2$. The sufficient additional count is $2r(\eps^{-1}-1)$, and
\[
 k-(3r+1)-2r(\eps^{-1}-1)
 \ge 2r/\eps-r-1>r-1\ge0.
\]
Lemma~\ref{lem:scalar} completes the proof. Zero tails terminate exactly.
\end{proof}

\section{Other fixed summable power envelopes}
The preceding inverse-square theorem has a much smaller constant than the following general estimate. The latter is included to make the range of the method explicit.

\begin{theorem}[Tail-power envelope]\label{thm:power}
Fix $p>1$, let
\[
 c=\max\{2,\lceil(p-1)^{-1}\rceil\},\qquad
 d_p=p+1+(p-1)^{-1},\qquad K_p=(3d_p/p)^p,
\]
and define
\begin{align}
 B_p^0&=\frac{c}{c-1}\frac{K_p(c+1)^{1+1/c}}{c^p}
             \left(1+\frac{K_p}{c^p}\right)^{1/c},\notag\\
 C_p&=c+2+B_p^0. \label{eq:Cp}
\end{align}
If $\lambda_{r+j}(A)\le a j^{-p}$ for $1\le j\le n-r$, where $a=\lambda_{r+1}(A)>0$, then \eqref{eq:target} holds with $C=C_p$. For fixed $p$, this constant is independent of $n,r,A$, and $\eps$.
\end{theorem}
\begin{proof}
A generating-function coefficient estimate gives, for $z>0$,
\[
 e_\ell(1,2^{-p},3^{-p},\ldots)
 \le z^{-\ell}\prod_{j=1}^\infty(1+zj^{-p}).
\]
The logarithm of the product is at most
\[
 \int_0^\infty\log(1+zx^{-p})\,dx
 =z^{1/p}\int_0^\infty\log(1+u^{-p})\,du
 <d_p z^{1/p}.
\]
Indeed, on $(0,1)$ the integrand is bounded by $\log2-p\log u$, and on $(1,\infty)$ it is bounded by $u^{-p}$. Minimize $z^{-\ell}\exp(d_pz^{1/p})$ at $z^{1/p}=p\ell/d_p$ and use $\mathrm e<3$. The result is
\begin{equation}
 e_\ell^{\rm tail}\le a^\ell(K_p/\ell^p)^\ell. \label{eq:powercoeff}
\end{equation}
This is the usual generating-function estimate, here applied to the tail rather than the full spectrum.

For $m\ge1$, combine \eqref{eq:head}, \eqref{eq:split}, and \eqref{eq:powercoeff} to obtain
\[
 e_{r+m}(A)\le P_r a^m(K_p/m^p)^m(1+K_p/m^p)^r.
\]
Proposition~\ref{prop:path} yields
\begin{equation}
 \frac{f_{r+m}}{\tau_r(A)}\le
 \frac{m}{m-1}((r+m)!)^{1/m}\frac{K_p}{m^p}
 (1+K_p/m^p)^{r/m}\quad(m>1). \label{eq:powermean}
\end{equation}
Take $m=cr$. Using $(r+m)!\le(r+m)^{r+m}$ and $1+1/c-p\le0$ gives
\[
 \frac{f_{(c+1)r}}{\tau_r(A)}\le
 \frac{c}{c-1}\frac{K_p(c+1)^{1+1/c}}{c^p}
 (1+K_p/c^p)^{1/c}r^{1+1/c-p}\le B_p^0.
\]
Use the larger warm-start constant $B=1+B_p^0>1$. With $C_p=c+1+B$, the count $\lceil C_pr/\eps\rceil$ supplies $(c+1)r$ starting pivots and at least $rB/\eps$ more. Lemma~\ref{lem:scalar} proves the result, with termination handling the dimension cap and a zero tail.
\end{proof}

\begin{remark}[The tail assumption is not automatic]
The index $j$ in these hypotheses begins at the first eigenvalue \emph{after} the target rank. For example, the assumption $\lambda_{r+j}\le\lambda_{r+1}j^{-p}$ demands a factor $2^{-p}$ decrease already at $j=2$. A global sequence $\lambda_i=i^{-p}$ need not satisfy that tail-normalized assumption at a large target rank. Nor does an arbitrary finite summable spectrum satisfy it with a fixed $p>1$. The constants above cannot be made universal by silently choosing an input-dependent envelope or exponent.
\end{remark}

\section{A unit tail-rate clock comparison is impossible}
Attach to each nonzero residual an independent exponential holding time of rate $T$. At its end use the pivot law \eqref{eq:update}. Equivalently, the transition to $R^{(j)}$ has rate $R_{jj}$. The embedded sequence is exactly RPCholesky. Let $N(t)$ be the number of positive pivots by time $t$.

\begin{theorem}[No affine mean-count bound with unit tail rate]\label{thm:clock}
For every finite real $a\ge0$, there is a finite real positive-definite matrix with $\tau_1(A)=1$ and a time $t>0$ such that
\begin{equation}
 \E N(t)>a+t. \label{eq:clockfailure}
\end{equation}
Thus no finite universal $a$ makes $\E N(t)\le ar+t\tau_r(A)$ valid on all inputs. The same obstruction applies to any tail-rate coefficient $b\le1$.
\end{theorem}
\begin{proof}
For an integer $n\ge2$ set
\begin{equation}
 \delta=\frac1{n-1},\qquad A=\mathbf1\mathbf1^*+\delta I_n.
 \label{eq:clockmatrix}
\end{equation}
Its eigenvalues are $n+\delta$ and $\delta$ repeated $n-1$ times, so $\tau_1(A)=1$. The Schur complement after any $k<n$ distinct pivots has the unselected block
\[
 \delta I_{n-k}+\frac{\delta}{\delta+k}\mathbf1\mathbf1^*.
\]
Consequently the holding rates are deterministic:
\begin{equation}
 T_k=\delta(n-k)\left(1+\frac1{k+\delta}\right)\quad(0\le k<n),
 \qquad T_n=0. \label{eq:spikerate}
\end{equation}
Every rate after the first pivot is at most $2$. The total count is therefore bounded by one initial pivot plus a rate-two candidate Poisson count; in particular, with $m(t)=\E N(t)$,
\begin{equation}
 m(t)\le1+2t. \label{eq:clockupper}
\end{equation}

For every integer $0\le k\le n$,
\begin{equation}
 T_k\ge1+\frac1{k+1}-\frac{k}{n-1}. \label{eq:clocklower}
\end{equation}
For $1\le k<n$, use $\delta\le1$ and $(n-k)/(n-1)=1-(k-1)/(n-1)$; the subtracted coefficient is
$(k-1)(1+1/(k+1))=k-2/(k+1)\le k$. At $k=0$ the assertion is immediate, and at $k=n$ its right side is negative.

The finite-state generator gives $m'(t)=\E T_{N(t)}$. Apply \eqref{eq:clocklower}, Jensen's inequality for $x\mapsto1/(x+1)$, and \eqref{eq:clockupper}:
\[
 m'(t)\ge1+\frac1{m(t)+1}-\frac{m(t)}{n-1}
 \ge1+\frac1{2+2t}-\frac{1+2t}{n-1}.
\]
Integration from $m(0)=0$ proves the explicit estimate
\begin{equation}
 \boxed{\displaystyle
 \E N(t)\ge t+\frac12\log(1+t)-\frac{t+t^2}{n-1}.}
 \label{eq:clockbound}
\end{equation}

Now choose finite integer parameters
\begin{equation}
 J=\lceil4(a+2)\rceil,\qquad t=2^J,\qquad n=t^2+t+2. \label{eq:clockparams}
\end{equation}
The last fraction in \eqref{eq:clockbound} is less than one. Also
$\log2=\int_1^2x^{-1}\,dx>1/2$, so
$\frac12\log(1+t)>J/4\ge a+2$. Thus $\E N(t)>t+a+1$, proving \eqref{eq:clockfailure}. All matrices in this proof are finite. For $b\le1$, the proposed upper bound $a+bt$ is no larger than $a+t$.
\end{proof}

For example, $a=1$ gives $J=12$, $t=4096$, and $n=16\,781\,314$. The exact scalar certificate in the archive proves the relevant finite inequalities without allocating that matrix. This is an analytic existence certificate, not a sampled numerical mean.

The third package rules out a leading rank coefficient exactly one even when its tail-rate coefficient is an arbitrary finite positive constant. An additive rank coefficient below one also fails on exact-rank inputs, whose eventual count equals their rank while their spectral tail is zero. Together, these observations show that an affine mean-count comparison of this form would need \emph{both} coefficients strictly greater than one. They do not rule out such comparisons with larger coefficients. Nor do they disprove RA-01: the rank-one matrices \eqref{eq:clockmatrix} satisfy the stronger discrete bound in Lemma~\ref{lem:small}.

\section{Why the determinant comparison does not finish RA-01}
\begin{proposition}[Saturation on diffuse tails]\label{prop:saturation}
Fix positive integers $r,m$, and let $k=r+m$. There are diagonal positive-definite inputs for which the coefficient $M_{r,m}(A)$ in \eqref{eq:M} approaches $k!/m!$ arbitrarily closely. Consequently, for a fixed integer $c\ge1$, the integrated bound \eqref{eq:integrated} at $k=(c+1)r$ is not bounded by an absolute multiple of $\tau_r$ as $r$ grows.
\end{proposition}
\begin{proof}
For $N\ge m$ and $L>1$, take
\[
 A_{L,N}=\operatorname{diag}
 \left(L^r,L^{r-1},\ldots,L,
 \underbrace{1/N,\ldots,1/N}_{N\text{ entries}}\right).
\]
Its tail $\tau_r$ is one. Let $P_r=L^{r(r+1)/2}$. For fixed $r,m,N$, each spectral tail $\tau_i$ with $i<r$ is asymptotic to its largest eigenvalue as $L\to\infty$, so $D_r/P_r\to1$. In the head--tail expansion \eqref{eq:split}, every term omitting a head factor vanishes after division by $P_r$. Hence
\[
 \lim_{L\to\infty}M_{r,m}(A_{L,N})
 =k!\binom Nm N^{-m}.
\]
Letting $N\to\infty$ gives $k!/m!$. Both limits can be approximated with finite integer parameters.

For $m=cr$ the corresponding root obeys
\[
 \left(\frac{((c+1)r)!}{(cr)!}\right)^{1/(cr)}
 =\left(\prod_{j=1}^{r}(cr+j)\right)^{1/(cr)}
 \ge(cr)^{1/c}\longrightarrow\infty.
\]
The additional factor $m/(m-1)$ in \eqref{eq:integrated} cannot remove this divergence.
\end{proof}

\textbf{This is a limitation of the displayed upper comparison, not a lower bound on RPCholesky's error.} These are diagonal matrices, and the first package proves the desired diagonal guarantee with $C=2$. In particular, the saturation calculation cannot be used as a counterexample to the algorithm. It explains exactly why the current tail-coefficient argument needs extra tail structure.

The remaining unrestricted assertion is still the existence of absolute $a\ge1$ and $B>1$ such that
\begin{equation}
 f_{\min\{n,\lceil ar\rceil\}}(A)\le B\tau_r(A)
 \quad\text{for every allowed }A,r. \label{eq:gap}
\end{equation}
Equation \eqref{eq:target} implies \eqref{eq:gap} by taking $\eps=1/2$. Conversely, \eqref{eq:gap} and Lemma~\ref{lem:scalar} imply \eqref{eq:target}, for example with $C=a+B+1$. Theorems~\ref{thm:geo}, \ref{thm:square}, and \ref{thm:power} prove \eqref{eq:gap} on their stated classes only.

The inverse-arrival-time argument from the third package remains a possible route. A stochastic comparison
\[
 N(t)\le_{\rm st}\lceil ar\rceil+\operatorname{Poisson}(b\tau_r(A)t)
\]
with absolute coefficients would suffice there. It is not established here. A first-moment count inequality cannot simply be substituted for stochastic domination of all upper tails. No general dense-matrix argument closing \eqref{eq:gap} has been obtained in this package.

\section{Verification, reproduction, and audit limits}
The standard-library program \texttt{code/verify\_exact.py} uses only rational arithmetic. Its recorded execution passes \textbf{8,500 explicitly counted checks}, organized in the JSON report by purpose. The checks cover eight rational matrix cases, alternate-order agreement of Schur residuals, actual RP pivot probabilities, one-step trace identities, the retained determinant bound and its integrated form, and the tail coefficient estimates. They also cover the rank-two scalar chain, factorial constants for ranks 3 through 100, 90 rank--tolerance pairs, and finite scalar clock and saturation certificates.

The rational matrices have known spectra because they are diagonal or rational Householder conjugates of specified nonnegative diagonal spectra. The program is not a general PSD-certification system. Some small matrices reach the dimension cap before the warm-start count; the JSON explicitly records this. The nontrivial finite probability and coefficient checks remain distinct from a universal theorem.

The optional NumPy program \texttt{code/verify\_floating.py} exhausts the subset distributions of two order-11 matrices, one real and one genuinely complex Hermitian. It also runs 1,200 Monte Carlo trajectories across six additional real or complex matrices. These tests exercise non-capped warm-start counts. They are floating-point diagnostics, not rational certificates, and Monte Carlo means are not used in any proof.

From the archive root:
\begin{verbatim}
python code/verify_manifest.py
python code/verify_exact.py
# Optional: NumPy required.
OPENBLAS_NUM_THREADS=1 python code/verify_floating.py
# Optional: a LaTeX installation is required.
sh paper/build.sh
\end{verbatim}

Run the manifest check before rebuilding outputs: timestamps or regenerated timing fields can change bytes without changing mathematical content. The preceding round-three ZIP is retained unchanged in \texttt{provenance/}; it in turn contains the earlier packages. The copied exact pivot engine is attributed to the first package. Source dependencies and the distinction between earlier and current claims are recorded in \texttt{SOURCE\_AUDIT.md}.

\textbf{Audit limit.} No independent human or separate-agent mathematical review, and no Lean, Coq, Isabelle, or other proof-assistant verification was performed. Exact arithmetic verifies the displayed finite calculations, not the quantifiers over all dimensions, the analytic generating-function proofs, or the stochastic arguments. Those claims rest on the written proofs and remain subject to mathematical review. The archive is a partial-results package, not a complete solution of RA-01. No repository files were modified.

\clearpage
\begin{thebibliography}{9}
\bibitem{repo} Open Problems in Numerical Linear Algebra. \emph{RA-01---Optimal pivot count for RPCholesky trace approximation}. Canonical repository entry, accessed September 13, 2026.\newline
\url{https://github.com/ajt60gaibb/OpenProblemsInNLA/blob/main/randomized-and-low-rank-approximation/RA-01/README.md}.

\bibitem{epperly} Ethan N. W. Epperly. \emph{A new analysis of the randomly pivoted Cholesky algorithm}. arXiv:2608.20633v1, posted August 21, 2026. See Section 2, Corollary 1.3, and the conjecture after (1.5).\newline
\url{https://arxiv.org/html/2608.20633v1}.

\bibitem{chen} Yifan Chen, Ethan N. Epperly, Joel A. Tropp, and Robert J. Webber. \emph{Randomly pivoted Cholesky: Practical approximation of a kernel matrix with few entry evaluations}. Communications on Pure and Applied Mathematics 78(5) (2025), 995--1041. See Algorithm 1 and Lemmas 5.3--5.5.\newline
\url{https://doi.org/10.1002/cpa.22234}; accessible manuscript:
\url{https://arxiv.org/html/2207.06503v6}.

\bibitem{divan} Ryan Divan and Marc Aur\`ele Gilles. \emph{Convergence rates of randomly pivoted methods for low-rank approximation}. arXiv:2609.06287v1, September 5, 2026. See Section 3 and Appendix A.\newline
\url{https://arxiv.org/html/2609.06287v1}.

\bibitem{dlmf} NIST Digital Library of Mathematical Functions. \emph{Section 4.36, Infinite Products and Partial Fractions}, equation 4.36.1. Accessed September 13, 2026.\newline
\url{https://dlmf.nist.gov/4.36.E1}.

\bibitem{prior} Prior conversation proof packages. \emph{RA-01: Sharp rank-one bounds, diagonal lower bounds, and block-structured guarantees}; \emph{RA-01: Sharp column-selection lower bounds and dense-matrix guarantees}; and \emph{RA-01: Arrival-time bounds and limits of clock comparisons}. Unpublished research notes, not independently reviewed. Preserved through \texttt{provenance/RA01\_round3\_clock\_analysis.zip}.
\end{thebibliography}
\end{document}
